\documentclass[11pt]{article}
\usepackage[margin=1in]{geometry}
\usepackage{amsmath,amssymb,amsthm}
\usepackage{hyperref}

\newtheorem{theorem}{Theorem}
\newtheorem{lemma}{Lemma}
\newtheorem{proposition}{Proposition}
\newtheorem{conjecture}{Conjecture}
\theoremstyle{definition}
\newtheorem{definition}{Definition}

\newcommand{\RR}{\mathbb{R}}
\DeclareMathOperator{\conv}{conv}

\title{Two hidden layers compute the maximum of six numbers}
\author{Matthew Sun}
\date{2026}

\begin{document}
\maketitle

\begin{abstract}
We prove that $\max(x_1,\dots,x_6)$ is exactly computable by a ReLU neural network with two hidden
layers of arbitrary width. This resolves an open problem of Bakaev, Brunck, Hertrich, Stade, and
Yehudayoff, who established the five-input case and explicitly left the six-input case open, and which is
also listed as open in subsequent work. We give an explicit construction as a signed combination of six
$S_6$-orbits of $P_2$ support functions with denominator $360$, together with an exact, floating-point
free computer-assisted proof of correctness. We then give computational evidence that six is the exact
threshold: $\max_7$ and $\max_8$ are not representable in the same family even after enlarging it, and we
reduce the question of a genuine two-versus-three hidden-layer separation to a single normal-form lemma.
\end{abstract}

\section{Introduction}

A feedforward ReLU network with $k$ hidden layers computes a continuous piecewise-linear (CPWL) function,
and conversely every CPWL function arises this way~\cite{ABMM}. A central question is how the number of
hidden layers limits expressivity, and the canonical test function is the maximum
$\max_n(x)=\max(x_1,\dots,x_n)$. A balanced binary tree of pairwise maxima computes $\max_n$ with
$\lceil\log_2 n\rceil$ hidden layers, and Hertrich, Basu, Di~Summa, and Skutella~\cite{HBDS} conjectured
this is optimal: that exactly representing $\max_n$ requires $\lceil\log_2(n+1)\rceil$ hidden layers, which
would make the ReLU depth hierarchy strict.

The conjecture holds for integer weights: Haase, Hertrich, and Loho~\cite{HHL} proved an integer-weight
lower bound of $\lceil\log_2 n\rceil$ hidden layers via lattice-polytope volumes. Over the rationals the
picture differs. Bakaev, Brunck, Hertrich, Stade, and Yehudayoff~\cite{BBHSY} showed that $\max_5$ needs
only two hidden layers, disproving the $n=5$ case and improving the general upper bound, and they left
open whether two hidden layers suffice for six inputs.

\begin{theorem}\label{thm:main}
$\max(x_1,\dots,x_6)$ is computable by a ReLU network with two hidden layers.
\end{theorem}

The construction uses rational weights, as it must: the integer-weight lower bound~\cite{HHL} would
otherwise force three layers. Our proof is computer-assisted but exact; every arithmetic step is over the
integers or rationals, and no floating-point computation enters the certificate.

\section{Preliminaries}

For a polytope $P\subset\RR^n$ its support function is $h_P(x)=\max_{g\in P}\langle g,x\rangle$. Then
$\max_n=h_\Delta$ for the standard simplex $\Delta=\conv\{e_1,\dots,e_n\}$. Support functions are additive
under Minkowski sum, $h_{P+Q}=h_P+h_Q$, and satisfy $h_{\conv(A\cup B)}=\max(h_A,h_B)$; two polytopes with
equal support functions are equal.

Following~\cite{BBHSY}, define $P_0=$ points and
$P_{k+1}=\{\sum_i P_i * Q_i : P_i,Q_i\in P_k\}$, where $*$ is join (convex hull of union) and the sum is
Minkowski. Then $P_1$ is the class of zonotopes. A function is computable by a $k$-hidden-layer ReLU
network if and only if it is a difference $h_Q-h_P$ with $P,Q\in P_k$~\cite{BBHSY,HHL}. Expanding the
Minkowski sums and using additivity, the two-hidden-layer functions are exactly the signed linear
combinations of building blocks
\[
   h_R(x)=\max_{g\in R}\langle g,x\rangle, \qquad R=\conv(Z_1\cup Z_2)\ \text{a join of two zonotopes.}
\]
To compute $\max_6$ in two layers it therefore suffices to write $h_\Delta$ as a signed sum of such $h_R$.

\paragraph{The weight-2 lattice.} We take building-block vertices on the weight-2 lattice points of the
dilated simplex $2\Delta^{5}$, namely $\{2e_i\}$ and $\{e_i+e_j\}$ (21 points in $\RR^6$); zonotopes are
generated by braid directions $e_a-e_b$ restricted so that all vertices stay on the lattice. This is the
lattice underlying the construction of~\cite{BBHSY} for $\max_5$.

\paragraph{Symmetry.} $\max_6$ is $S_6$-invariant, so averaging any representation over $S_6$ yields a
symmetric one; we therefore search over $S_6$-orbits of building blocks, collapsing $16390$ two-way joins
to $78$ orbit classes and making an exact linear solve tractable.

\section{The construction}

Solving exactly for orbit coefficients reproducing $\max_6$ and minimizing the number of terms yields a
construction with six orbit terms and denominator $D=360$:
\[
   \max_6 = \sum_{o=1}^{6} c_o \Big(\sum_{R\in o} h_R\Big), \qquad
   c_o \in \Big\{ \tfrac{1}{360}, -\tfrac{1}{360}, \tfrac{1}{360}, \tfrac{1}{90}, -\tfrac{1}{90}, -\tfrac{1}{180} \Big\},
\]
with zero linear part. The six orbit representatives, as vertex (gradient) sets, are listed in the
accompanying file \texttt{construction\_max6.txt}. Each $R$ is a join of two weight-2 zonotopes, so each
$h_R$ is a two-hidden-layer function and the signed sum is two hidden layers. Expanding the six orbits
over $S_6$ gives $1530$ individual $P_2$ terms.

\section{The proof of Theorem~\ref{thm:main}}

We prove the displayed identity holds for all $x\in\RR^6$. Both sides are CPWL, so equality is equivalent
to equality of gradients on every full-dimensional cell of their common arrangement.

\paragraph{Reduction to a small arrangement.} On any region the construction's gradient equals
$\sum_t c_t\,\arg\max(R_t)$, where $\arg\max(R_t)$ is the vertex of $R_t$ maximizing $\langle g,x\rangle$.
This selection changes only across the hyperplanes $g-g'=0$ for $g,g'$ vertices of a common term, the
weight-2 difference hyperplanes. By $S_6$-symmetry it suffices to work in the closed fundamental chamber
$x_1\ge\cdots\ge x_6$, where $\max_6=x_1$ and the target gradient is $e_1$. In gap coordinates
$y_k=x_k-x_{k+1}\ge 0$ a difference hyperplane with normal $h$ becomes
$\sum_k\big(\sum_{i\le k} h_i\big)y_k=0$, which is \emph{forced} (constant sign on $y>0$) unless its
prefix sums change sign. Of the $120$ distinct difference hyperplanes, $85$ are forced; only $35$ are
unforced, giving an arrangement of $35$ hyperplanes in $\RR^5$.

\paragraph{(1) Complete cell enumeration.} We enumerate every full-dimensional cell of this arrangement in
the open cone $y>0$ by a deterministic depth-first search over sign vectors, pruning infeasible branches.
There are exactly $2608$ cells, and for each we produce an exact integer interior point.

\paragraph{(2) Exact gradient.} At each cell's integer witness we compute, in exact integer arithmetic,
the gradient $\sum_t c_t\,\arg\max(R_t)$. On all $2608$ cells it equals $e_1$. At an interior point every
relevant comparison is strict, so each $\arg\max$ is unique and the gradient is well defined and constant
on the cell.

\paragraph{(3) Completeness by adjacency closure.} It remains to show the $2608$ cells are all of them,
without relying on floating-point pruning. The regions of a hyperplane arrangement inside the connected
cone $y>0$ form a connected graph under facet adjacency, and adjacent regions differ in exactly one sign.
For every enumerated cell and every single-sign flip, the neighbor is either another enumerated cell or
empty. Emptiness is certified exactly: by Gordan's theorem the open cell $\{y:\langle g,y\rangle>0\ \forall
g\}$ is empty iff there exist nonnegative rationals $x$, not all zero, with $\sum_j x_j g_j=0$ over the
generators (the orthant generators and the signed normals); we exhibit such an $x$ and verify the identity
over the rationals. Across all $2608\times 35$ flips, every neighbor is in-set ($13756$) or exactly
certified empty ($77524$), with none uncertified. A nonempty subset of a connected graph closed under
adjacency is the whole graph, so the enumeration is complete.

By (1)--(3) the construction's gradient equals $e_1$ on every cell of the chamber, so the construction
equals $x_1=\max_6$ throughout the chamber; by continuity and $S_6$-symmetry it equals $\max_6$ on all of
$\RR^6$. No floating-point arithmetic appears in (2) or (3). \qed

The certificate is reproduced by \texttt{prove\_max6.py}; an independent CPWL verifier agrees.

\section{The threshold and a route to a separation}

The two-hidden-layer functions whose building blocks have vertices in the weight-2 lattice form a finite
family: a weight-2 zonotope has vertices among the finitely many lattice points, and the building blocks
are the two-way joins of such zonotopes (including degenerate point zonotopes, giving pyramids). We
enumerate this family completely (all zonotopes generated by arbitrary, not merely braid, weight-2 point
differences) and solve over $\mathbb{Q}$.

\begin{proposition}\label{prop:max7}
$\max_7$ is not a signed sum of support functions of two-way joins of weight-2 zonotopes. Equivalently,
$\max_7$ has no two-hidden-layer representation with all building-block vertices in the weight-2 lattice.
\end{proposition}

This is exact: the complete family yields a linear system over $\mathbb{Q}$ that is inconsistent for
$\max_7$ (and consistent for $\max_6$, recovering Theorem~\ref{thm:main}). The same holds for $\max_8$.
This locates a qualitative change exactly at $n=6$, and we conjecture:

\begin{conjecture}
$\max_n$ is computable in two hidden layers if and only if $n\le 6$.
\end{conjecture}

Proposition~\ref{prop:max7} is a genuine lower bound over the weight-2 lattice; what is missing for an
unconditional separation is only that the lattice may be assumed weight-2. The reduction $\max_n$ two-layer $\iff \Delta+P=Q$ with $P,Q\in P_2$ is standard; one checks
that any representation may be symmetrized, and that $\Delta+P=Q$ forces the normal fan of $Q$ to refine
the braid fan. The remaining content is a normal-form lemma:

\begin{conjecture}[Lattice normal form]
If $\max_n$ is computable in two hidden layers, it has a representation with all building-block vertices in
the weight-2 lattice.
\end{conjecture}

Granting this, two-layer representability becomes a finite computation, and the $n=7$ negative becomes the
statement that $\max_7$ requires three hidden layers, a separation of two from three hidden layers over the
reals. We note that the conformity route is closed: the optimal shallow construction is provably not
braid-conforming (indeed our $\max_6$ construction requires non-conforming blocks), consistent with the
conditional bounds of Grillo, Hertrich, and Loho~\cite{GHL}. The relevant tools for the lattice normal
form are the decomposition polyhedra of Brandenburg, Grillo, and Hertrich~\cite{BGH} and deformation-cone
theory for the braid fan. Contemporary lower bounds are consistent with Theorem~\ref{thm:main}: \cite{HHL}
is integer-weight only, \cite{AHM} is for finite-precision weights, \cite{GHL} is for braid-conforming
networks, and \cite{Saf} is a width rather than depth bound at depth at least three.

\begin{thebibliography}{9}
\bibitem{ABMM} R.~Arora, A.~Basu, P.~Mianjy, A.~Mukherjee. Understanding Deep Neural Networks with
Rectified Linear Units. ICLR 2018. arXiv:1611.01491.
\bibitem{HBDS} C.~Hertrich, A.~Basu, M.~Di~Summa, M.~Skutella. Towards Lower Bounds on the Depth of ReLU
Neural Networks. NeurIPS 2021; SIAM J. Discrete Math. 37(2):997--1029, 2023. arXiv:2105.14835.
\bibitem{HHL} C.~Haase, C.~Hertrich, G.~Loho. Lower Bounds on the Depth of Integral ReLU Neural Networks
via Lattice Polytopes. ICLR 2023. arXiv:2302.12553.
\bibitem{BBHSY} E.~Bakaev, F.~Brunck, C.~Hertrich, J.~Stade, A.~Yehudayoff. Better Neural Network
Expressivity: Subdividing the Simplex. STOC 2026. arXiv:2505.14338.
\bibitem{BGH} M.~Brandenburg, M.~Grillo, C.~Hertrich. Decomposition Polyhedra of Piecewise Linear
Functions. ICLR 2025. arXiv:2410.04907.
\bibitem{GHL} M.~Grillo, C.~Hertrich, G.~Loho. Depth-Bounds for Neural Networks via the Braid Arrangement.
NeurIPS 2025. arXiv:2502.09324.
\bibitem{AHM} G.~Averkov, C.~Hojny, M.~Merkert. On the Expressiveness of Rational ReLU Neural Networks
With Bounded Depth. ICLR 2025. arXiv:2502.06283.
\bibitem{Saf} I.~Safran. A Depth Hierarchy for Computing the Maximum in ReLU Networks via Extremal Graph
Theory. 2026. arXiv:2601.01417.
\end{thebibliography}

\end{document}
